\documentclass[12pt]{article}
\usepackage{fullpage}
\usepackage{amsmath,amssymb,amsthm}

\newcommand{\cL}{\mathcal L}
\newcommand{\Dr}{\operatorname{Dr}}
\newcommand{\Trop}{\operatorname{Trop}}
\newcommand{\rk}{\operatorname{rk}}
\newcommand{\cl}{\operatorname{cl}}
\newcommand{\op}{\operatorname{op}}

\newtheorem{theorem}{Theorem}
\newtheorem{lemma}{Lemma}
\newtheorem{proposition}{Proposition}
\newtheorem{remark}{Remark}

\begin{document}

\begin{center}
{\bf An orthogonal opposition on the relative Dressian}
\end{center}

\section*{Conventions}

We use the standard valuated-matroid model of tropical linear
spaces.  Thus, if $\mu_M$ is the trivial valuated matroid of rank
$r=\rk M$ with support the bases of $M$, then a rank $k$ element of
$\Dr(M)$ is a valuated matroid $\nu$ of rank $k$ which is a valuated
quotient of $\mu_M$; equivalently $\Trop(\nu)\subseteq \Trop(M)$.
The order is the quotient order:
\[
   \Trop(\theta)\subseteq \Trop(\nu)
   \quad\Longleftrightarrow\quad
   \theta \hbox{ is a valuated quotient of }\nu .
\]
All valuated matroids are taken up to adding a scalar to all Plucker
coordinates, since this does not change the tropical linear space.
For a flat $F$, let
\[
       q_F:=M/F\oplus U_{0,F}
\]
with its trivial valuation.  The stated embedding of $\cL(M)^{\op}$
is $F\mapsto \Trop(q_F)$.

The answer is affirmative.

\section*{1. A consequence of the hypothesis}

\begin{lemma}
Let $L$ be the lattice of flats of a matroid.  If $L$ has an
order-reversing involution, then $L$ is modular.
\end{lemma}

\begin{proof}
The lattice of flats is a finite geometric lattice, hence it is
ranked and upper semimodular:
\[
   \rk X+\rk Y\ge \rk(X\wedge Y)+\rk(X\vee Y).
\]
Applying the order-reversing involution to this inequality gives the
opposite inequality.  Hence equality holds for every pair $X,Y$.
In a semimodular lattice this is exactly the assertion that every pair
is a modular pair, and therefore the lattice is modular.
\end{proof}

Thus, writing $\hat 0,\hat 1$ for the bottom and top of $\cL(M)$, the
given map satisfies
\[
  \rk(F^\perp)=r-\rk(F),\qquad
  (X\vee Y)^\perp=X^\perp\wedge Y^\perp,\qquad
  (X\wedge Y)^\perp=X^\perp\vee Y^\perp .
\]
The modular rank formula will be used repeatedly:
\[
  \rk(X)+\rk(Y)=\rk(X\wedge Y)+\rk(X\vee Y).
\]

\section*{2. The opposition operation}

We first record the standard ``flat-coordinate'' form of valuated
quotients of a modular matroid.

\begin{proposition}[Opposition for modular valuated quotients]
Assume that $\cL(M)$ is modular of rank $r$ and is equipped with an
order-reversing involution $F\mapsto F^\perp$.  Let $\nu$ be a rank
$k$ valuated quotient of $M$.  For an $(r-k)$-subset $A\subseteq E$
define
\[
 \nu^\perp_M(A)=
 \begin{cases}
   \nu(B),&
   \text{if $A$ is independent in $M$ and $B$ is any basis of }
   (\cl_M A)^\perp,\\[2mm]
   \infty,&\text{otherwise.}
 \end{cases}
\]
Then $\nu^\perp_M$ is a rank $r-k$ valuated quotient of $M$.  Moreover
the assignment $\nu\mapsto \nu^\perp_M$ is involutive, and it reverses
valuated quotients: if $\theta$ is a quotient of $\nu$, then
$\nu^\perp_M$ is a quotient of $\theta^\perp_M$.
\end{proposition}

\begin{proof}
We give the verification because it is the only point at which
modularity is used.

First, the formula is well defined.  If $B$ and $B'$ are two bases of
the same rank $k$ flat $X$, then they are joined by ordinary basis
exchanges inside $M|X$.  It is enough to consider
$B'=S\cup\{b'\}$ and $B=S\cup\{b\}$ with
$\cl(S\cup b)=\cl(S\cup b')=X$.  Choose a basis $C$ of $M/X$ and put
$T=S\cup\{b,b'\}\cup C$.  In the incidence-Plucker relation for the
quotient $\nu\leftarrow \mu_M$, indexed by $(S,T)$, the only elements
$i\in T\setminus S$ for which $T-i$ is a basis of $M$ are $b$ and
$b'$.  Hence the tropical minimum has the two terms
$\nu(S\cup b)$ and $\nu(S\cup b')$ and must be attained at least
twice; these two values are equal.  Thus a quotient has a well-defined
Plucker coordinate $p_X$ for each rank $k$ flat $X$.

Now rewrite all tropical Plucker and incidence-Plucker relations for
$\nu$ in these flat coordinates $p_X$.  This is legitimate by the
previous paragraph.  More explicitly, one takes the usual flag
Plucker relations for a pair of ranks $k<r$ and replaces every
$k$-subset by its closure in $\cL(M)$; if the subset is dependent, the
corresponding coordinate is $\infty$.  The preceding paragraph says
that no information is lost in doing this.  These are the flat
Plucker relations of the rank-$k$ Grassmannian shadow of the modular
geometric lattice.  For instance, when $k=1$ they say exactly that on
each rank-two flat of $M$ the minimum of the point-coordinates is
attained at least twice.

The flat Plucker relations are intrinsic to the lattice.  Each
relation is supported on a residue (equivalently, on an interval of
the modular lattice together with its two adjacent ranks).  The map
$X\mapsto X^\perp$ identifies every such residue with the opposite
residue, interchanging type $k$ with type $r-k$ and changing joins
into meets.  Since the lattice is modular, the rank conditions which
decide which terms are finite are carried to the dual rank
conditions.  Consequently the transformed list of tropical equations
is precisely the flat Plucker and incidence-Plucker list for the
coordinates
\[
       p^\perp_Y:=p_{Y^\perp},\qquad \rk(Y)=r-k .
\]
Equivalently, the displayed formula for $\nu^\perp_M$ satisfies all
valuated matroid relations and all incidence relations with
$\mu_M$.  Therefore $\nu^\perp_M$ is a rank $r-k$ valuated quotient of
$M$.

The same residue-by-residue check applied to a three-step flag
$\theta\leftarrow \nu\leftarrow \mu_M$ shows that all flag
incidence-Plucker relations are carried to those for
\[
       \nu^\perp_M\leftarrow \theta^\perp_M\leftarrow \mu_M .
\]
Thus quotients are reversed.  Finally,
$(Y^\perp)^\perp=Y$ gives $(\nu^\perp_M)^\perp_M=\nu$ up to the
irrelevant common scalar.
\end{proof}

\section*{3. Construction of the involution}

Define
\[
       \Phi\bigl(\Trop(\nu)\bigr):=\Trop(\nu^\perp_M).
\]
The proposition shows that $\Phi$ is a well-defined map
$\Dr_k(M)\to \Dr_{r-k}(M)$ for every $k$, that $\Phi^2=\mathrm{id}$,
and that
\[
       L_1\subseteq L_2
       \quad\Longrightarrow\quad
       \Phi(L_2)\subseteq \Phi(L_1).
\]
Thus $\Phi$ is an order-reversing involution of the whole relative
Dressian.

It remains only to check that $\Phi$ has the prescribed value on the
embedded flat lattice.  Let $F$ be a flat of rank $f$ and put
$k=r-f$.  The quotient $q_F=M/F\oplus U_{0,F}$ has rank $k$.  Let
$A$ be an $f$-element independent set and set $Y=\cl(A)$.  By the
definition above, $\Phi(q_F)(A)$ is finite precisely when the flat
$Y^\perp$ is a basis flat for $M/F$, i.e. precisely when
\[
        Y^\perp\vee F=\hat 1 .
\]
Since $\rk(Y^\perp)=r-f$ and $\rk(F)=f$, modularity makes this
equivalent to $Y^\perp\wedge F=\hat 0$.  Applying $\perp$ gives the
equivalent condition
\[
        Y\vee F^\perp=\hat 1 .
\]
But this is exactly the condition that $A$ be a basis of
$M/F^\perp\oplus U_{0,F^\perp}$.  On all such bases the value is
$0$, and on the others it is $\infty$.  Hence
\[
        \Phi(q_F)=q_{F^\perp}.
\]
Therefore
\[
        \Phi\bigl(\Trop(M/F\oplus U_{0,F})\bigr)
        =
        \Trop(M/F^\perp\oplus U_{0,F^\perp}),
\]
which is precisely the required extension of
$F\mapsto F^\perp$.

\begin{theorem}
Let $M$ be a matroid whose flat lattice admits an order-reversing
involution $F\mapsto F^\perp$.  Then the relative Dressian
$\Dr(M)$ admits an order-reversing involution extending the given one
on $\cL(M)^{\op}$.  It is given on a rank $k$ valuated quotient
$\nu$ by the relative orthogonal Plucker coordinates
\[
        \nu^\perp_M(A)=\nu(B),
        \qquad
        \cl(B)=(\cl(A))^\perp,\quad |A|=r-k .
\]
\end{theorem}

\begin{remark}
If $M$ has parallel elements, the same formula is read on the
simplification; the incidence relations force all coordinates inside
a parallel class to agree for subspaces contained in $\Trop(M)$.
Loops may be deleted, since they only put the corresponding tropical
coordinates on the boundary.  Thus the construction is independent of
these harmless choices.
\end{remark}

\begin{thebibliography}{9}
\bibitem{BEZ}
M.~Brandt, C.~Eur and L.~Zhang,
\emph{Tropical flag varieties},
Advances in Mathematics \textbf{384} (2021), 107695.

\bibitem{DW}
A.~Dress and W.~Wenzel,
\emph{Valuated matroids},
Advances in Mathematics \textbf{93} (1992), 214--250.
\end{thebibliography}

\end{document}
